\section{RL for LLM Reasoning}

\subsection{超越偏好：可验证奖励}

在数学、编程等领域，答案是否正确可以\textbf{自动验证}。此时不需要人类偏好，可以直接用正确性作为奖励进行 RL 训练。

\begin{importantbox}{RL for Reasoning 的兴起}
DeepSeek-R1 等模型展示了 RL 训练能让 LLM 学会推理中的自我纠错和思维链等能力。这是 RL 在 LLM 中除对齐之外的另一个重要应用方向。
\end{importantbox}

\subsection{可验证奖励的设计}

可验证任务常见信号包括：
\begin{itemize}
    \item \textbf{逻辑一致性}：自动判别数学解答是否满足步骤逻辑。
    \item \textbf{合成推理}：用 symbolic executor 检查推理链条，reward 直接反馈到 token。
    \item \textbf{生成式验证}：用 verifier LM 对 candidate 回答打分，作为 auxiliary reward。
\end{itemize}

\begin{table}[H]
\centering
\begin{tabular}{p{0.3\textwidth} p{0.35\textwidth} p{0.32\textwidth}}
\toprule
任务类别 & 可验证信号 & 典型指标 \\
\midrule
数学推理 & steps correctness & chain-of-thought accuracy、full solution correctness \\
编程 & unit test pass rate & pass@k、runtime exceptions \\
知识问答 & factuality consistency & TruthfulQA、F1 against evidence \\
\bottomrule
\end{tabular}
\caption{可验证奖励覆盖的任务与指标}
\end{table}

\begin{knowledgebox}{自监督 verifier 的作用}
当人工比较不可行时，引入自监督 verifier（如 program executor 或 verifier LM）可以生成 reward signal，保障 reasoning RL 的闭环。
\end{knowledgebox}

\subsection{实验与评估}

DeepSeek-R1 等工作常用 baseline：Instruction tuned LM + RLHF contrast、DPO、RL for reasoning pipeline。验证方法包括 LM 生成、多轮提示和 external verifier 检查。

\begin{importantbox}{Reasoning RL 的洞察}
\textit{``RL 为模型提供了自我纠错的机制，而不是仅仅加强模仿''}——这种反馈让 model 在面对长链推理或 multi-hop question 时更愿意 traverse deeper reasoning tree。
\end{importantbox}

评估时也要考虑 inference latency，因为 reasoning RL 常常需要 rollouts 和 verifier，部署时需要权衡 real-time 体验。

\subsection{本章小结}

RL 在 LLM 中不仅用于对齐（RLHF/DPO），也服务于 reasoning：通过可验证奖励与 verifier pipeline，模型得到自我纠错的监督；实验中需要同时考量 correctness 与 latency trade-off。

